\documentclass[12pt,a4paper]{report}

\usepackage[utf8]{inputenc}
\usepackage[T1]{fontenc}
\usepackage[spanish]{babel}

\title{LabMate\\Memoria del proyecto}
\author{Alex, Oscar y Aritz}
\date{\today}

\begin{document}

\maketitle

\chapter{Introducción, problema y objetivos}

\section{Nombre y frase de venta}

El proyecto se llama \textbf{LabMate}.

Nuestra frase de venta es: ``Tu propia IA, segura, agéntica y rápida: vive en tu ordenador, conoce tu proyecto y trabaja a tu lado''.

\section{Problema o necesidad y usuario}

Hoy mucha gente usa IAs de fuera para trabajar en sus proyectos. Pero esto tiene varios problemas:

\begin{itemize}
  \item Se pierde tiempo copiando y pegando código.
  \item Estas IAs no conocen nuestro proyecto.
  \item A veces se inventan cosas.
  \item Los datos salen de nuestro ordenador.
\end{itemize}

Pensamos en estudiantes y en desarrolladores que quieren ayuda sin perder tiempo ni privacidad.

\section{Objetivos generales}

Con LabMate queremos:

\begin{itemize}
  \item Crear un asistente de IA que funcione en local y sea privado.
  \item Que se use desde la terminal, dentro de la carpeta de un proyecto.
  \item Que responda sobre lo que hay en esa carpeta, resuma documentos y guarde notas.
  \item Que diga ``no lo sé'' en lugar de inventar.
  \item Que al principio solo lea, y que nunca vea contraseñas ni claves.
\end{itemize}

\end{document}
